\documentclass[book.tex]{subfiles}

\begin{document}
\chapter{Neural Radiance Fields}

\label{chap:nerfs}
\noindent\rule{11cm}{0.4pt} \\
\textbf{Prerequisites:} \autoref{chap:molecular_hamiltonian}, \autoref{chap:dft}, \autoref{chap:hartree_fock} \\
\textbf{Difficulty Level:} ** \\
\noindent\rule{11cm}{0.4pt}
\vspace{5mm}

Neural radiance fields (NeRFs) provide a new methodology for transforming input coordinates into volume densities and radiances. The core mathematical transformation is given by

\begin{align}
(x, y, z, \theta, \phi) \mapsto (R,G, B, \sigma)
\end{align}

Here $x,y,z$ are spatial coordinates, and $(\theta, \phi)$ are viewing angles. This transformation is typically performed by a fully connected network $f_\theta$. The NeRF provides dramatic improvements in view synthesis, allowing for a global view to be learned from a collection of input images with known orientations
\begin{align}
\mathcal{D} &= \{ x_i\} \\
x_i &\in \mathbb{R}^{m \times n}
\end{align}
The core idea is that rays of light can be sent backwards from the image data to compute a representation of the radiance at any given point. This radiance can be used to project out new sampled images from new orientations.


\end{document}